% !TEX root = ../main.tex

% !TEX root = ../main.tex

\section{Further Protocol Details}
\label{app:zq-protocol}

\subsection{The Clearing Rule}
\label{app:zq-clearing-rule}

We define \(\mathsf{Clear}\) for books containing ordinary and imbalance-only (IO) orders. A book has \(N\) real price ticks, padded to \(n\) rows, and consists of four nonnegative integer sequences \(b_i,a_i,b_i^{io},a_i^{io}\), with zero entries for \(i\ge N\). A market buy is recorded at tick \(N-1\) and a market sell at tick 0. For each tick, let
\[
D_i=\sum_{j\ge i}b_j,\quad S_i=\sum_{j\le i}a_j,\quad
V_i=\min(D_i,S_i),\quad g_i=S_i-D_i,\quad \delta_i=|g_i|.
\]
Similarly, \(D_i^{io}=\sum_{j\ge i}b_j^{io}\) and \(S_i^{io}=\sum_{j\le i}a_j^{io}\). We call demand \emph{long} if \(g_i<0\), supply \emph{long} if \(g_i>0\), and the book \emph{balanced} if \(g_i=0\). At an imbalanced tick, the at-price long-side volume is \(\lambda_i=b_i\) when demand is long and \(\lambda_i=a_i\) when supply is long.

\paragraph{Clearing receipt.} Let \(V=\max_i V_i\). If \(V=0\), the outcome is \(\textsf{no-cross}\). Otherwise, define
\begin{itemize}
\item the maximizing plateau \(\Pi=\{i:V_i=V\}\), with \(c=\min\Pi\) and \(d=\max\Pi\);
\item the minimum imbalance \(v=\min_{i\in\Pi}\delta_i\);
\item a candidate tick \(i\in\Pi\) satisfying \(\delta_i=v\) and either \(v=0\) or \(\lambda_i\ge v\).
\end{itemize}
Choose \(p^*\) as the lowest candidate. The resulting receipt is \((V,c,d,p^*,v,s,P,J)\), where \(s=P=J=0\) if \(v=0\). Otherwise, \(s=0\) when demand is long and \(s=1\) when supply is long, \(P=\lambda_{p^*}\), and \(J=S_{p^*}^{io}\) if \(s=0\) or \(J=D_{p^*}^{io}\) if \(s=1\). The condition \(\lambda_{p^*}\ge v\) ensures that unfilled ordinary volume can be rationed at the clearing price without reducing fills of better-priced orders. In the main paper's notation, \(V=V_{max}\) and \(v=V_{min\Delta}\).

\subsubsection{Fills}
\label{app:zq-fills}

The exchange admits \(I=\min(J,v)\) eligible IO units, leaving residual imbalance \(R=v-I\). The total traded volume is \(V+I\). For an order of quantity \(q\), the unrounded execution entitlement is
\begin{itemize}
\item \(q\), for an ordinary order strictly better priced than \(p^*\), or on the short side at \(p^*\);
\item \(q(P-R)/P\), for an ordinary order at \(p^*\) on the long side;
\item \(qI/J\), for an eligible IO order on the short side;
\item zero, for any other order.
\end{itemize}
The ratios are used only when their denominators are positive. If \(v=0\), all eligible ordinary orders fill in full and IO orders receive zero. Integer rounding is left to the venue.

\subsection{Encoding and Public Polynomials}

Let \(\mathbb F\) be the scalar field of BN254 and \(H=\{1,\omega,\ldots,\omega^{n-1}\}\) an evaluation domain, where \(n\) is a power of two, \(N\le n\), and \(\omega\) has order \(n\). The four book columns \(B,A,IOB,IOA\) are polynomials represented by their evaluations on \(H\), with \(B(\omega^i)=b_i\), for example. Each column entry is below \(2^{k_{in}}\), and accumulated quantities are bounded by \(2^k\), with \(k\ge k_{in}+\lceil\log_2 N\rceil\). Our 1,000-tick, 32-bit configuration uses \(n=1024\) and \(k=42\).

The nine derived value columns consist of four depth accumulators for ordinary and IO orders; executable volume \(Min\); two surplus columns; imbalance \(\Delta\); and the inequality witness \(Slack\). Together with the four book columns, these give 13 value columns. For the configuration above, the range checks require 46 limb columns and the lookup argument 29 more columns, for 88 witness columns in total.

Let \(L_i(X)\) denote the Lagrange polynomial selecting tick \(i\). The verifier constructs the public polynomials
\[
\begin{aligned}
Pad(X)&=\sum_{i=N}^{n-1}L_i(X),\\
Mask_P(X)&=\sum_{i=c}^{d}L_i(X),\\
Mask_C(X)&=\sum_{i\in\{c-1,d+1\}\cap\{0,\ldots,n-1\}}L_i(X),\\
Mask_T(X)&=\begin{cases}
L_{p^*-1}(X),&p^*>c\text{ and }(s=1\text{ or }v=0),\\
0,&\text{otherwise.}
\end{cases}
\end{aligned}
\]
The range table is likewise a public polynomial with evaluations \(T(\omega^i)=i\) on the table rows (0 through 255 for eight-bit limbs). No commitments to these public polynomials are needed.

\subsection{Clearing Constraints}

\paragraph{Scalar checks.} Before checking polynomial identities, the verifier checks the receipt:
\[
\begin{aligned}
0&\le c\le p^*\le d\le N-1, & 1&\le V<2^k,\\
s&\in\{0,1\}, & 0&\le v\le P<2^k,\quad 0\le J<2^k.
\end{aligned}
\]
If \(v=0\), it also requires \(s=P=J=0\). The polynomial constraints F1--F33 are listed in the main paper; here we explain only the less immediate parts.

\paragraph{Maximum-volume plateau.} F14 bounds the volume at each tick by \(V\), F15--F16 pin the claimed endpoints to that maximum, and F17 makes the immediately adjacent ticks strictly smaller:
\[
Mask_C(X)\bigl(V-Min(X)-1-Slack(X)\bigr)=0.\tag{F17}
\]
These conditions identify the entire plateau because \(Min\) is the pointwise minimum of a non-increasing demand sequence and a non-decreasing supply sequence. Specifically, for \(i<j<\ell\), \(Min[j]\ge\min(Min[i],Min[\ell])\). The selector \(Mask_C\) omits a cliff at a domain boundary, and the range check on \(Slack\) enforces strict inequality without publishing the cliff volumes.

\paragraph{Price tie-break.} F19--F20 establish minimum imbalance on the plateau and pin the chosen side. The verifier also checks \(v\le P\). For supply-long or balanced outcomes with \(p^*>c\), F21 excludes a lower qualifying price:
\[
Mask_T(X)\bigl(\Delta(X)-v-1-Slack(X)\bigr)=0.\tag{F21}
\]
The same \(Slack\) column supports the inequalities in F17 and F21.

\paragraph{No-cross.} When ordinary orders cannot cross, the exchange proves a smaller statement using the ordinary book columns, their cumulative depths, F1--F8, the relevant range checks, and
\[
AccA(X)\,AccB(X)=0.
\]
Since both depths are nonnegative, this forces zero executable volume at every tick. IO orders do not create a cross.

\subsection{Range Checks and Lookup}

To prove that a column \(Y\) represents integers in \([0,2^b)\), we decompose it into \(w\)-bit limb columns \(\ell_0,\ldots,\ell_{m-1}\), with \(m=\lceil b/w\rceil\), and impose
\[
Y(X)-\sum_{j=0}^{m-1}2^{wj}\ell_j(X)=0.
\]
Each limb must belong to the public range table. We use \(w=8\) when \(n\ge256\), and narrower limbs for smaller domains. If \(b\) is not divisible by \(w\), we also check that \(2^{w-r}\ell_{m-1}\) lies in the table, for \(r=b\bmod w\); this restricts the top limb to \(r\) bits.

For table membership we use a logarithmic-derivative lookup. Let \(f_1,\ldots,f_E\) be the looked-up columns and \(T\) the table column. The prover commits to a multiplicity column \(Mult\); after challenge \(\beta\), it supplies pair helpers \(h_j\), a table helper \(U\), and a running-sum column \(S\), satisfying
\[
\begin{aligned}
h_j&=\frac{1}{\beta-f_{2j-1}}+\frac{1}{\beta-f_{2j}},&
U&=\frac{Mult}{\beta-T},\\
S(\omega X)-S(X)&=\sum_j h_j(X)-U(X).
\end{aligned}
\]
The corresponding polynomial constraints are
\[
\begin{aligned}
h_j(\beta-f_{2j-1})(\beta-f_{2j})
 &-(\beta-f_{2j-1})-(\beta-f_{2j})=0,\\
U(\beta-T)-Mult&=0,\\
S(\omega X)-S(X)-\sum_j h_j(X)+U(X)&=0.
\end{aligned}\tag{F23}
\]
Summing the last identity over the cyclic domain cancels the running-sum terms and gives equality between the rational lookup sums and table multiplicities. Since the looked-up columns and multiplicities are committed before \(\beta\), an invalid membership claim is detected with high probability. The running sum needs no initial condition.

\subsection{Proving the Constraints with a PIOP}
\label{app:piop}

We follow a PLONK-style polynomial interactive oracle proof \cite{GWC19}, with an additional lookup round. The main steps are:
\begin{enumerate}
\item Commit to the value columns, range-check limbs, and \(Mult\). The four book-column commitments already appear on the public board.
\item After challenge \(\beta\), commit to the lookup-helper columns.
\item After challenge \(\alpha\), combine constraints into \(F(X)=\sum_t\alpha^tF_t(X)\), commit to the quotient \(Q(X)=F(X)/Z_H(X)\), where \(Z_H(X)=X^n-1\).
\item After challenge \(\zeta\), open the columns at \(\zeta\). Shifted openings at \(\omega\zeta\) are also required for the four depth accumulators, \(A\), \(IOA\), and \(S\).
\item Reconstruct \(F(\zeta)\) from the openings and check \(F(\zeta)=Q(\zeta)Z_H(\zeta)\).
\end{enumerate}
We use hiding KZG commitments \cite{KZG10} over BN254, aggregate openings at each point, and apply Fiat--Shamir using a SHA-256 transcript bound to the parameters, receipt, and commitments. Random multiples of \(Z_H\) mask witness columns without changing their values on \(H\). For the 1,024-row, 32-bit configuration, the clearing proof occupies 6,124 bytes; the reference string supports degrees through \(2n+8\). The public selector evaluations can depend on plateau width, even though the number of commitments and opening points is fixed.

\subsection{Committed Book and Trader Checks}
\label{app:committed-book}

The clearing proof is tied to the final commitments of an append-only order log. Trader receipts establish inclusion of acknowledged orders; an additional update-validity proof prevents logged negative updates from removing them.

\subsubsection{Commitment and Opening}

Let \([x]_1\) denote a scalar multiple of the first pairing-group generator, and let \(\Gamma=[\gamma]_1\) denote the blinding generator. A book column is committed as
\[
C=[\mathrm{col}(\tau)]_1+[\rho_0(\tau)Z_H(\tau)]_1+[\gamma r(\tau)]_1,
\qquad \mathrm{col}(X)=\sum_i\mathrm{col}[i]L_i(X),
\]
where \(\rho_0\) and \(r\) are random degree-three masking polynomials. Their values are retained for the clearing proof. To certify an initially empty column, the exchange publishes \(C_0\), \(C_\rho=[\rho_0(\tau)]_1\), and \(C_r=[\gamma r(\tau)]_1\), with a Schnorr proof \cite{Sch91} of knowledge of the blinding polynomial, and the verifier checks
\[
e(C_0-C_r,[1]_2)=e(C_\rho,[Z_H(\tau)]_2).
\]
Thus the committed value component is a multiple of \(Z_H\) and vanishes at every tick.

\subsubsection{Recording an Order}

For an order of quantity \(q\) at tick \(\ell\), the exchange updates the selected column commitment by \(\Phi=q[L_\ell(\tau)]_1+\rho\Gamma\), while reblinding the other three columns. It appends all four updated commitments to the public board in a fixed order and returns the trader the receipt
\[
(\mathrm{column},\ell,q,\rho,C_{before},C_{after}).
\]
The trader verifies
\[
C_{after}=C_{before}+q[L_\ell(\tau)]_1+\rho\Gamma
\]
and checks that the two commitments are consecutive entries for that column in the public log. The differences are Pedersen-style commitments \cite{Ped91a}; fresh blinding conceals which column carries the actual order. Each log entry hashes the preceding digest, its index, column, and commitment. The inclusion check relies on participants seeing a consistent published log head.

\subsubsection{Update Validity}
\label{app:zq-update-validity}

A valid receipt alone would not stop the exchange from cancelling the order with a subsequent negative update. Let \(\Phi_0,\ldots,\Phi_{\mu-1}\) be the logged commitment differences, four per order. The update-validity proof establishes that each step adds a quantity \(q_j\in[0,2^{k_{in}})\) at one valid tick \(\ell_j\), or adds zero.

Use a second domain \(K\) of size \(M\), with \(K\supseteq H\), and one row per logged step. The prover commits to \(Qty[j]=q_j\) and \(Tick[j]=\omega^{\ell_j}\), along with an empty masking commitment \(\Phi_{mask}\), certified as above. After challenge \(\eta\), both sides form the aggregate commitment \(C_W\) and its underlying polynomial evaluations
\[
C_W=\Phi_{mask}+\sum_{j<\mu}\eta^j\Phi_j,
\qquad w[i]=\sum_{j:\ell_j=i}\eta^j q_j.
\]
The verifier computes \(C_W\) from the log, while the prover supplies the committed polynomial \(W(X)\) with evaluations \(w[i]\). A weighted lookup, with challenge \(\beta\) drawn after commitments, establishes
\[
\sum_{j<\mu}\frac{\eta^j q_j}{\beta-Tick[j]}
=\sum_{i<n}\frac{w[i]}{\beta-\omega^i}.
\]
This is enforced by four groups of polynomial constraints. Let \(\kappa\) generate \(K\), with \(\kappa^{M/n}=\omega\). The public polynomial \(P_\eta\) takes value \(\eta^j\) on step rows \(j<\mu\) and zero elsewhere; on row \(iM/n\), \(X_H=\omega^i\) and \(S_H=1\), with zero selector values on other rows. The prover also supplies helper columns \(G,U_H,S_2\). Write \(m_{in}=\lceil k_{in}/w\rceil\). The constraints are:
\[
\begin{aligned}
Qty(X)-\sum_{j<m_{in}}2^{wj}\ell_j^{Qty}(X)&=0, &&\text{(U1)}\\
G(X)(\beta-Tick(X))-P_\eta(X)Qty(X)&=0, &&\text{(U2)}\\
U_H(X)(\beta-X_H(X))-S_H(X)W(X)&=0, &&\text{(U3)}\\
S_2(\kappa X)-S_2(X)-G(X)+U_H(X)&=0. &&\text{(U4)}
\end{aligned}
\]
The quantity limbs in U1 also undergo the table lookup. Cyclic summation of U4 gives the weighted lookup identity. Because \(\eta\) is sampled after the logged steps are fixed, a false aggregate would require a coincidental cancellation. Zero-quantity steps permit the three blinding-only updates per order. For 32-bit quantities, this proof has 17 commitments and occupies 1,700 bytes. With one row per logged step, the domain accommodates up to \(M/4\) orders.

